%----------
%	OBJECTIVES AND MOTIVATION
%----------	

\color{red}
[Toda la historia de \#seacabó]
\color{black}


The 20th of August 2023 was the day set for the 2023 FIFA Women's World Cup final, between Spain and England. It was the first time in history that either team had reached that stage, so it was a transcendent event that had all eyes on. Spain ended up winning that match 1-0 with the goal of Real Madrid player Olga Carmona. Unfortunately, that day ended up with a bittersweet flavor because of the "Rubiales Affair". At the moment of the medals ceremony, former RFEF (Real Federación Española de Fútbol) president Luis Rubiales kissed Spanish player Jenni Hermoso on the lips, with a visible discomfort from her side.

The RFEF called an extraordinary general meeting for 25th August and was expected the resignation of Rubiales. However, he 

The social movement against the conduct of former RFEF president Luis Rubiales and the structured women denigration in the sport is widely considered as the Spanish \textit{#MeToo} movement. It was rapidly called \textit{"\#seacabó"} due to a tweet\footnote{From \url{https://x.com/alexiaputellas/status/1695052143978295707?lang=en}, last visited 02/07/24} that Spanish player and Ballon d'Or winner Alexia Putellas published after the RFEF extraordinary general meeting, and that was later accompanied by teammates and other relevant sports figures.


\begin{figure}[H]
	\ffigbox[\FBwidth]
	{\caption{Tweet published by Alexia Putellas in support of teammate Jenni Hermoso}\label{tweet_seacabo}} 
 % [Aquí ponemos como queremos que aparezca en el índice, sin la referencia]{Aquí como queremos que aparezca en el documento, con la referencia}
	{\includegraphics[scale=0.75]{Plantilla_TFG_ingles_2019/imagenes/aputellas_tweet.png}}
\end{figure}



The Spanish society, including politicians of all political backgrounds and men's football players and teams, condemned Rubiales in an unprecedented display of unity that showed support for women's football. Rubiales ended up resigning as president on September 10, 2023 and was subsequently banned from football activities for three years. There is currently an investigation into allegations of coercion in a criminal case against him.

Therefore, the aim of this work is to aspire to contribute to the creation of safer online environments, where women can express themselves freely without any fear of harm or abuse. While AI could be a strong ally in controlling online content, it is essential that social media platforms provide more reliable mechanisms that protect free speech while also keeping people, especially women, safe.

\color{red}
Consequently, relevant available datasets are limited in size, highly imbalanced, and exhibit topical and lexical biases. https://www.taylorfrancis.com/books/mono/10.1201/9781032654829/regulating-hate-speech-created-generative-ai?refId=985b6d7e-7548-4820-af81-4f4845b339e1&context=ubx
\color{black}

Therefore, we could identify the following objectives for this thesis:

\textbf{General Objective:} The general objective of this thesis is to address the detection of hate speech on social media platforms in the Spanish language through the use of automatic methods.

\textbf{Specific Objective 1:} Our first specific objective is to conduct a comprehensive study of the State of the Art in hate speech detection, focusing on the comparison between Machine Learning, Deep Learning, and Generative AI models.

\textbf{Specific Objective 2:} Our second specific objective is to implement traditional text encoding models, such as TF-IDF, and Bag of Words, and explore more complex embeddings like FastText, Word2Vec, BERT, and RoBERTa, to enhance the performance of HS detection models.

\textbf{Specific Objective 3:} Our third specific objective is to develop and apply various data preprocessing techniques to improve the accuracy and efficiency of HS detection models. This includes tasks such as tokenization, stopword removal, and handling imbalanced data through techniques like SMOTE and ADASYN.

\textbf{Specific Objective 4:} Our fourth specific objective is to design and develop a range of Machine Learning models, Deep Learning models, and Generative AI models for the detection of hate speech in Spanish social media content.

\textbf{Specific Objective 5:} Our fifth specific objective is to systematically evaluate the performance of the implemented models and embeddings, comparing their metrics and performance in the context of HS detection.

Therefore, the motivation of this work is to create reliable and unbiased automatic detection models that help identify hate speech towards women in social networks \textit{(X/Twitter)}. In particular, by means of artificial intelligence (AI) and Natural Language Processing (NLP) algorithms.

\newpage